\documentclass[11pt]{article}

\usepackage{amsmath,amssymb,amsthm,mathtools}
\usepackage[T1]{fontenc}
\usepackage{lmodern}
\usepackage[margin=1in]{geometry}
\usepackage{enumitem}
\usepackage{hyperref}
\usepackage{microtype}

\hypersetup{
  colorlinks=true,
  linkcolor=blue,
  citecolor=blue,
  urlcolor=blue
}

\newtheorem{theorem}{Theorem}[section]
\newtheorem{proposition}[theorem]{Proposition}
\newtheorem{corollary}[theorem]{Corollary}
\newtheorem{definition}[theorem]{Definition}
\newtheorem{lemma}[theorem]{Lemma}
\newtheorem{remark}[theorem]{Remark}

\newcommand{\Z}{\mathbb Z}
\newcommand{\Acal}{\mathcal A}
\newcommand{\Pad}{\mathsf{Pad}}
\newcommand{\Ev}{\mathsf E}
\newcommand{\one}[1]{\mathbf 1_{\{#1\}}}

\title{One-pass odd-$m$ globalization package}
\author{}
\date{March 2026}

\begin{document}
\maketitle

\begin{abstract}
This note rewrites the odd-$m$ D5 globalization chain as one read-through
theorem package.

The package is self-contained in the following theorem-level sense:
\begin{enumerate}[label=\textup{(\arabic*)},leftmargin=2.5em]
\item every theorem used later is stated in this document;
\item every downstream proof after the standing package inputs is carried out
      in this document;
\item no cross-note citation is needed while reading the argument.
\end{enumerate}

The standing package inputs are kept explicit and honest:
\begin{enumerate}[label=\textup{(\arabic*)},leftmargin=2.5em]
\item the exact trigger-family formula for $H_{L1}$;
\item the coarse patch facts needed for pre-exit patch avoidance;
\item the componentwise concrete bridge package on full chains;
\item the regular chart / endpoint support package;
\item the exceptional source-$3$ interior continuation fact.
\end{enumerate}

From those inputs, the note proves in one place the structural first-exit
theorem, the fixed-$\delta$ tail-length reduction, the chart/interface landing
theorem, the regular closure theorem, the actual exceptional interface landing
theorem, and the final odd-$m$ D5 globalization theorem.
\end{abstract}

\section*{Scope note}

This note is \emph{not yet} a graph-theoretic proof of Hamilton decomposition
for the odd-$m$, $d=5$ torus.

What it proves is the odd-$m$ D5 globalization theorem inside the accepted D5
regime and the fixed best-seed channel package.

To turn this document into a full Hamilton-decomposition proof, one still
needs:
\begin{enumerate}[label=\textup{(\arabic*)},leftmargin=2.5em]
\item a reduction theorem showing why this channel covers the whole odd-$m$,
      $d=5$ problem, or why it is the last unresolved channel;
\item a graph-level bridge from raw global $(\beta,\delta)$ exactness to a
      global edge-selection rule, global $2$-factors, Hamiltonicity of each
      factor, and edge-disjoint coverage of the full torus edge set;
\item closure of the standing Inputs~A--E and an explicit quantifier for the
      exact allowed values of $m$.
\end{enumerate}

So the right description of this note is: a one-document theorem package for
the odd-$m$ D5 globalization theorem, not yet a standalone proof of Hamilton
decomposition for the full torus.

\section{Scope and basic definitions}

Fix an odd modulus $m$ in the accepted D5 regime and the best-seed channel
\[
R1 \to H_{L1}
\]
with source seed words
\[
L=[2,2,1], \qquad R=[1,4,4].
\]
Write
\[
R_{\mathrm{reg}}=\{1,2,4,5,\dots,m-1\}
\]
for the regular source labels.

\begin{definition}[full chains and actual lifts]
A \emph{full chain} is a complete boundary block between consecutive visits to
the $\Theta=2$ section.

An \emph{actual lift} of a realized boundary label $\delta_0$ is an actual
accessible full-chain occurrence with boundary label $\delta_0$.
\end{definition}

\begin{definition}[tail length]
The \emph{remaining tail length} $\tau(C)$ of an actual lift $C$ is the number
of full-chain successors available after $C$ before terminality. If successors
continue forever, then $\tau(C)=\infty$.
\end{definition}

\begin{definition}[endpoint class and right-congruence]
Fix the standard accepted terminal padding convention on future current-event
words. Two actual lifts lie in the same \emph{endpoint class} when they
determine the same padded future current-event word.

Write $\Pad$ for the common terminal padding tail supplied by that convention.
The corresponding boundary right-congruence class is denoted by $\rho$.
\end{definition}

\begin{definition}[distinguished rows]
The \emph{exceptional cutoff row} is
\[
\delta=3m-3.
\]
The corresponding interface rows are
\[
3m-2
\qquad\text{and}\qquad
3m-1.
\]
\end{definition}

\begin{definition}[accessible boundary union and splice-connected components]
The \emph{accessible boundary union} is the union of all actual accessible full
chains.

A \emph{splice-connected accessible component} is a connected component of that
union under the successor / predecessor splice relation between consecutive full
chains.
\end{definition}

\begin{definition}[regular and exceptional realizations]
A full chain or actual lift is \emph{regular} when it lies in one of the
regular source labels $u\in R_{\mathrm{reg}}$.

A full chain or actual lift is \emph{exceptional} when it lies in the
exceptional source label $u=3$.

A \emph{regular realization} of a realized boundary label $\delta$ is a regular
actual lift at that label.
\end{definition}

\section{Standing theorem-package inputs}

Everything after this section is proved in this note from the inputs listed
below.

\begin{proposition}[Input A: exact trigger family]
The target hole family for the unresolved best-seed channel is exactly
\[
H_{L1}=\{(q,w,u,\Theta)=(m-1,m-1,u,2):u\neq 2\}.
\]
Equivalently:
\begin{enumerate}[label=\textup{(\arabic*)},leftmargin=2.5em]
\item the target layer is $\Theta=2$;
\item the target current coordinates satisfy $q=w=m-1$;
\item the only excluded fiber is $u=2$.
\end{enumerate}
\end{proposition}

\begin{proposition}[Input B: coarse patch facts]
For the patched current classes, only the following coarse facts are used:
\begin{enumerate}[label=\textup{(\arabic*)},leftmargin=2.5em]
\item the classes $R1,R2,R3,L2$ require $w=0$;
\item the class $L3$ requires $w=1$;
\item a pre-exit occurrence of $L1$ would require the branch to have already
      reached the $w=m-1$ edge of the corridor.
\end{enumerate}
\end{proposition}

\begin{proposition}[Input C: componentwise concrete bridge]
On the unique $\Theta=2$ boundary state $(q,w,u,2)$ of a full chain, define
\[
\sigma = w+u-q-1 \pmod m,
\qquad
\delta=q+m\sigma.
\]
Then on each splice-connected accessible component:
\begin{enumerate}[label=\textup{(\arabic*)},leftmargin=2.5em]
\item consecutive full chains satisfy the odometer splice law
      \[
      \delta'=\delta+1 \pmod{m^2};
      \]
\item the current event class $\epsilon_4$ is determined by $(\beta,\delta)$;
\item the realized boundary image is a forward interval in $\Z/m^2\Z$ or the
      full odometer cycle.
\end{enumerate}
\end{proposition}

Let
\[
\Acal=
\{\mathrm{wrap},\mathrm{carry\_jump},\mathrm{other}_{1000},
  \mathrm{other}_{0010},\mathrm{flat}\}
\]
be the current-event alphabet, and write
\[
\Ev:\Z/m\Z\times \Z/m^2\Z\longrightarrow \Acal,
\qquad
(\beta,\delta)\longmapsto \epsilon_4(\beta,\delta)
\]
for the corresponding explicit event law.

\begin{proposition}[Input D: regular chart and endpoint support]
For each regular source label $u\in R_{\mathrm{reg}}$:
\begin{enumerate}[label=\textup{(\arabic*)},leftmargin=2.5em]
\item the $\Theta=2$ source window starts at
      \[
      s_u=m((u+2)\bmod m)-1;
      \]
\item that source window has exactly $m(m-3)$ consecutive full-chain labels;
\item its terminal label is
      \[
      e_u=s_u+m(m-3)-1=m(u-1)-2 \pmod{m^2};
      \]
\item grouped successor is exact inside the source window and advances labels
      by $\delta\mapsto \delta+1$;
\item every actual regular lift at the endpoint label $e_u$ continues on the
      true larger regular accessible union to the successor label
      \[
      e_u+1=s_{u-3}.
      \]
\end{enumerate}
\end{proposition}

\begin{definition}[regular source windows and regular holders]
For a regular source label $u\in R_{\mathrm{reg}}$, write
\[
I_u=[s_u,e_u]
\]
for the corresponding regular source window from Input~D.

Say that a boundary label $\delta$ has a \emph{source-$u$ regular holder} when
$\delta\in I_u$.

A regular holder is \emph{interior} when $\delta\neq e_u$, and \emph{terminal}
when $\delta=e_u$.
\end{definition}

\begin{proposition}[Input E: exceptional source-$3$ interior continuation]
Inside the exceptional source-$3$ chart, every exceptional full chain with
label strictly before the cutoff row $3m-3$ is an interior chart occurrence and
therefore continues internally by the same grouped successor rule
\[
\delta\mapsto \delta+1.
\]
\end{proposition}

\section{Structural first-exit theorem}

Use current coordinates $(q,w,u,\Theta)$ along the active branch after the
alternate-$2$ entry out of $R1$.

\begin{definition}[candidate orbit]
Define the source residue
\[
\rho:=u_{\mathrm{source}}+1 \pmod m,
\]
where $u_{\mathrm{source}}$ is the source-$u$ label of the chosen $R1$ family.
Start at the alternate-$2$ entry
\[
E=(m-1,1,u_{\mathrm{source}},2).
\]
Define the candidate orbit $\widehat x_n=(q_n,w_n,u_n,\Theta_n)$ by
\[
\widehat F(q,w,u,0)=(q+1,w,u,1),
\]
\[
\widehat F(q,w,u,1)=(q,w,u+1,2),
\]
\[
\widehat F(q,w,u,2)=\bigl(q,w+\one{q=m-1},u,3\bigr),
\]
\[
\widehat F(q,w,u,\Theta)=(q,w,u,\Theta+1)
\qquad (3\le \Theta\le m-1).
\]
\end{definition}

On the section $\Theta=2$, this induces the odometer return map
\[
(q,w)\mapsto (q+1,w+\one{q=m-1}).
\]
Equivalently, with
\[
\eta(q,w):=q+m(w-1)-(m-1),
\]
one has $\eta\mapsto \eta+1$.

The only two phase-$1$ states whose direction-$2$ or direction-$1$ successor can
reach $(q,w)=(m-1,m-1)$ are
\[
A=(m-1,m-2,1),
\qquad
B=(m-2,m-1,1),
\]
with indices
\[
\eta(A)=m(m-3),
\qquad
\eta(B)=m(m-2)-1.
\]
So $A$ is encountered before $B$.

\begin{lemma}[phase-$1$ source-residue invariant]
Along the candidate orbit, every phase-$1$ current state satisfies
\[
q\equiv u-\rho+1 \pmod m.
\]
Hence
\[
u(A)\equiv \rho-2 \pmod m,
\qquad
u(B)\equiv \rho-3 \pmod m.
\]
\end{lemma}

\begin{proof}
The first phase-$1$ state after entry is $(0,2,u_{\mathrm{source}},1)$, so the
relation holds there because $u_{\mathrm{source}}=\rho-1$.

From one phase-$1$ state to the next, the phase-$1$ witness step increments $u$
by $1$, the later phase-$0$ witness step increments $q$ by $1$, and no other
step changes either coordinate. So the congruence is preserved.
\end{proof}

\begin{lemma}[phase-$1$ trigger criterion]
Let $x$ be a phase-$1$ current state on the candidate orbit. An alternate step
from $x$ lands in $H_{L1}$ if and only if one of the following holds:
\begin{enumerate}[label=\textup{(\arabic*)},leftmargin=2.5em]
\item $x=A=(m-1,m-2,1)$ and $u(x)\neq 2$, in which case direction $2$ lands in
      $H_{L1}$;
\item $x=B=(m-2,m-1,1)$ and $u(x)\neq 2$, in which case direction $1$ lands in
      $H_{L1}$.
\end{enumerate}
\end{lemma}

\begin{proof}
By Input~A, the successor must satisfy
\[
(q',w',\Theta')=(m-1,m-1,2)
\qquad\text{and}\qquad
u'\neq 2.
\]
Only phase-$1$ states can step to layer $2$, and only directions $1$ and $2$
change $q$ or $w$. So the only two ways to reach $(q',w')=(m-1,m-1)$ are:
\begin{itemize}[leftmargin=2em]
\item direction $2$ from $(q,w)=(m-1,m-2)$;
\item direction $1$ from $(q,w)=(m-2,m-1)$.
\end{itemize}
Those are exactly $A$ and $B$. Directions $1$ and $2$ preserve $u$, so the only
extra condition is $u\neq 2$.
\end{proof}

\begin{proposition}[candidate first exits]
On the candidate orbit:
\begin{enumerate}[label=\textup{(\arabic*)},leftmargin=2.5em]
\item if $\rho\neq 4$ (regular family), the first exit to $H_{L1}$ occurs at
      $A=(m-1,m-2,1)$ by direction $2$;
\item if $\rho=4$ (exceptional family), the state $A$ does not exit and the
      first exit to $H_{L1}$ occurs at $B=(m-2,m-1,1)$ by direction $1$.
\end{enumerate}
Equivalently,
\[
T_{\mathrm{reg}}=(m-1,m-2,1),
\qquad
T_{\mathrm{exc}}=(m-2,m-1,1).
\]
\end{proposition}

\begin{proof}
By Lemma~3.3, only $A$ and $B$ can trigger an exit to $H_{L1}$. Since $A$ is
encountered first, it is enough to test whether $u(A)=2$.

By Lemma~3.2,
\[
u(A)\equiv \rho-2 \pmod m.
\]
So $u(A)=2$ if and only if $\rho=4$.

If $\rho\neq 4$, then $u(A)\neq 2$, so the first exit is at $A$ by direction
$2$. If $\rho=4$, then $u(A)=2$, so $A$ is blocked.

Still in the exceptional case, Lemma~3.2 gives
\[
u(B)\equiv \rho-3 = 1 \pmod m.
\]
So $u(B)\neq 2$, and the first exit is at $B$ by direction $1$.
\end{proof}

\begin{lemma}[pre-exit patch avoidance]
Let $N_f$ be the first-exit time of the candidate orbit in the chosen source
family $f$. For every $0\le n<N_f$, the candidate state $\widehat x_n$ avoids
all patched current classes.
\end{lemma}

\begin{proof}
The value $w$ starts at $1$ and changes only by the increment $w\mapsto w+1$ at
phase $2$ with $q=m-1$. So $w$ never decreases, and $w=0$ never occurs before
exit. By Input~B(1), the classes $R1,R2,R3,L2$ are impossible.

The value $w=1$ occurs only at the entry state $(m-1,1,u_{\mathrm{source}},2)$,
because the first phase-$2$ wrap immediately creates $w=2$ and thereafter $w$
never decreases. So by Input~B(2), the class $L3$ cannot occur after entry.
Since the entry state itself is the alt-$2$ start state rather than a patched
class, $L3$ is absent on the whole pre-exit segment.

On a regular family, Proposition~3.4 says the branch exits already at $A$,
before $w$ can ever reach $m-1$. So $L1$ is impossible.

On the exceptional family, the first creation of $w=m-1$ occurs at phase $3$
from $(m-1,m-2,2)$. The next phase-$1$ state then has $(q,w)=(0,m-1)$, not
$(m-1,m-1)$. The first later phase-$1$ state with $w=m-1$ that could also
carry $q=m-1$ would occur strictly after $B=(m-2,m-1,1)$, but
Proposition~3.4 says the branch exits already at $B$. So $L1$ is impossible
before exceptional exit as well.

Thus no patched current class is encountered before first exit.
\end{proof}

\begin{theorem}[actual first exits and pre-exit $B$-region invariance]
For each source family:
\begin{enumerate}[label=\textup{(\arabic*)},leftmargin=2.5em]
\item the actual active branch agrees with the candidate orbit up to first
      exit;
\item every pre-exit actual current state is $B$-labeled;
\item the first exit to $H_{L1}$ occurs at the universal family-dependent
      target $T_{\mathrm{reg}}$ or $T_{\mathrm{exc}}$, with regular families
      exiting at $T_{\mathrm{reg}}$ by direction $2$ and the exceptional family
      exiting at $T_{\mathrm{exc}}$ by direction $1$.
\end{enumerate}
\end{theorem}

\begin{proof}
Fix a source family and let $N_f$ be the candidate first-exit time from
Proposition~3.4. Prove by induction on $n\le N_f$ that the actual state $x_n$
agrees with the candidate state $\widehat x_n$ in the full current coordinates
$(q,w,u,\Theta)$.

At $n=0$, both states are the same entry state
\[
E=(m-1,1,u_{\mathrm{source}},2).
\]
By Lemma~3.5, it is not patched, hence it is $B$-labeled.

Now assume $x_n=\widehat x_n$ for some $n<N_f$. Lemma~3.5 says $\widehat x_n$
is not patched, so the actual state $x_n$ is $B$-labeled. Therefore its
witness update is the unmodified mixed-witness update, which induces exactly
the candidate step on $(q,w,u,\Theta)$: direction $1$ changes only $q$,
direction $4$ changes only $u$, direction $2$ changes only $w$, and the other
anchors leave $(q,w,u)$ fixed while increasing $\Theta$ by $1$. So
$x_{n+1}=\widehat x_{n+1}$ in full current coordinates.

Thus the actual and candidate branches agree for all $n\le N_f$, and every
pre-exit actual state is $B$-labeled.

For $n<N_f$, Proposition~3.4 and Lemma~3.3 show that the candidate state does
not exit to $H_{L1}$; because the trigger test depends only on the current
coordinates $(q,w,u,\Theta)$, the actual state does not exit either. At time
$N_f$, Proposition~3.4 gives an exit from the candidate orbit at the stated
universal target and direction, so the actual branch exits there as well.
\end{proof}

\section{Fixed-$\delta$ tail-length reduction}

\begin{definition}[one-chain event block]
For $\delta\in\Z/m^2\Z$, define the canonical one-chain event block
\[
W(\delta)=\bigl(\Ev(m-2,\delta),\Ev(m-3,\delta),\dots,\Ev(0,\delta),\Ev(m-1,\delta)\bigr).
\]
This is the current-event word read while the chain clock runs once through its
common boundary-to-boundary cycle.
\end{definition}

\begin{lemma}[the current chain word is determined by $\delta$]
Let $C$ be any actual lift with boundary label $\delta$. Then the current-event
word emitted along the chain $C$ is exactly $W(\delta)$.
\end{lemma}

\begin{proof}
By Input~C, every full chain starts at the common boundary clock value
$\beta=m-2$, and inside that chain the clock decreases deterministically by $1$
modulo $m$. The emitted current event at each position is therefore read from
the same explicit law $\Ev(\beta,\delta)$, with the same ordered list of clock
values $m-2,m-3,\dots,0,m-1$. So the full chain emits exactly the block
$W(\delta)$.
\end{proof}

\begin{proposition}[block stack determined by start label and tail length]
Let $C$ be an actual lift with boundary label $\delta$.
\begin{enumerate}[label=\textup{(\arabic*)},leftmargin=2.5em]
\item If $\tau(C)=t<\infty$, then the forward chain labels are
      \[
      \delta,\ \delta+1,\ \dots,\ \delta+t,
      \]
      where the chain at $\delta+t$ is terminal, and the padded future
      current-event word of $C$ is
      \[
      W(\delta)\,W(\delta+1)\,\cdots\,W(\delta+t)\,\Pad.
      \]
\item If $\tau(C)=\infty$, then the padded future current-event word of $C$ is
      the infinite concatenation
      \[
      W(\delta)\,W(\delta+1)\,W(\delta+2)\,\cdots.
      \]
\end{enumerate}
In particular, the padded future current-event word is completely determined by
$\delta$ together with $\tau(C)$.
\end{proposition}

\begin{proof}
Start with the current chain. By Lemma~4.2, it emits $W(\delta)$.

If no successor exists, then $\tau(C)=0$, the current chain is terminal, and by
the fixed padding convention the future word is $W(\delta)\Pad$. So the claim
holds in the base case.

Now suppose a successor chain exists. By Input~C, its boundary label is
$\delta+1$. By Lemma~4.2 again, that successor emits the block $W(\delta+1)$.
Repeating the same argument for each further successor yields the entire stack
$W(\delta),W(\delta+1),W(\delta+2),\dots$ until the first terminal chain, if
one occurs, or forever if no terminal chain occurs.

Thus a finite tail of length $t$ gives exactly the finite concatenation through
$\delta+t$ followed by the common padding suffix $\Pad$, while an infinite tail
gives the infinite concatenation. No other local future-law choice remains.
\end{proof}

\begin{theorem}[future data are determined by tail length]
Let $C_1$ and $C_2$ be actual lifts of the same realized boundary label
$\delta$. Then:
\begin{enumerate}[label=\textup{(\arabic*)},leftmargin=2.5em]
\item if $\tau(C_1)=\tau(C_2)$, then $C_1$ and $C_2$ have the same padded
      future current-event word;
\item $C_1$ and $C_2$ have the same padded future current-event word if and
      only if they lie in the same endpoint class;
\item $C_1$ and $C_2$ lie in the same endpoint class if and only if
      $\rho(C_1)=\rho(C_2)$.
\end{enumerate}
Equivalently, at fixed realized $\delta$, the padded future current-event word,
the endpoint class, and the boundary right-congruence class $\rho$ are all
determined by the remaining full-chain tail length.
\end{theorem}

\begin{proof}
By Proposition~4.3, the padded future current-event word of an actual lift is a
function only of its starting label $\delta$ and its tail length $\tau$. So
equality of tail length is sufficient for equality of the padded future word.

Items \textup{(2)} and \textup{(3)} are merely the definitions of endpoint
class and of the right-congruence label $\rho$ attached to the padded future
word.
\end{proof}

\begin{corollary}[different tail lengths create a later mixed-status row]
Let $C_1$ and $C_2$ be actual lifts of the same realized label $\delta$, and
assume $\tau(C_1)\neq\tau(C_2)$. Set
\[
t=\min\{\tau(C_1),\tau(C_2)\}<\infty.
\]
Then the later realized label $\delta+t$ is mixed-status: one actual lift at
that label is terminal and another actual lift at that same label continues.
\end{corollary}

\begin{proof}
By Proposition~4.3, both lifts follow the same deterministic block stack
through the chain labels $\delta,\delta+1,\dots,\delta+t$. At the last of these
labels, namely $\delta+t$, the lift with shorter tail has no further successor,
while the other lift still has at least one more successor. So $\delta+t$ is
realized with both a terminal lift and a continuing lift.
\end{proof}

\section{Regular chart consequences}

\begin{proposition}[chart/interface landing]
Read the actual exceptional post-exit continuation in the grouped full-chain
chart used for the boundary labels. Then its first two chart labels are forced
to be
\[
3m-2 \to 3m-1.
\]
Equivalently, the exceptional chain-label continuation is
\[
3m-3 \to 3m-2 \to 3m-1,
\]
where $3m-2$ is the unique terminal regular source-$4$ occurrence and $3m-1$
is the regular source-$1$ start.
\end{proposition}

\begin{proof}
The exceptional cutoff row is, by definition,
\[
\delta=3m-3.
\]
In the grouped chart quotient, successor is exact and advances the label by
$+1$. Therefore the first chart label after the cutoff row is forced to be
$3m-2$.

By Input~D, the source-$4$ terminal label is
\[
e_4=m(4-1)-2=3m-2,
\]
and the regular source-$1$ start label is
\[
s_1=m(1+2)-1=3m-1.
\]
Again by Input~D, the label immediately after $e_4$ is $s_{4-3}=s_1$. So the
interface is exactly $3m-2\to 3m-1$.
\end{proof}

\begin{theorem}[regular closure theorem]
On the true larger regular accessible union, every regular realization of every
realized $\delta$ continues.
\end{theorem}

\begin{proof}
Take any regular realization and choose one regular holder for its label.

If the realization is interior in that holder, grouped successor exactness
inside the source window gives a forward successor immediately.

If the realization sits at the terminal label of that holder, then its label is
some $e_u$, and Input~D gives a forward successor to $e_u+1=s_{u-3}$.

So every regular realization continues.
\end{proof}

\begin{corollary}[no regular mixed-status row]
No realized regular $\delta$ has both continuing and terminal regular lifts.
\end{corollary}

\begin{proof}
By Theorem~5.2, every regular lift continues.
\end{proof}

\section{Actual exceptional interface landing}

\begin{theorem}[actual exceptional interface landing]
Every actual lift of the exceptional cutoff row $\delta=3m-3$ reaches the
universal raw exit target
\[
T_{\mathrm{exc}}=(m-2,m-1,1)
\]
by direction $1$, and the corresponding full-chain labels continue as
\[
3m-3 \to 3m-2 \to 3m-1.
\]
Here $3m-2$ is the terminal regular source-$4$ occurrence and $3m-1$ is the
regular source-$1$ start.
\end{theorem}

\begin{proof}
By Theorem~3.6, every actual lift of the exceptional family exits the active
branch at the same raw state $T_{\mathrm{exc}}$ and with the same exit
direction $1$. Therefore all such lifts share the same immediate raw post-exit
continuation by raw determinism.

Proposition~5.1 then identifies the corresponding chain-label continuation with
\[
3m-3 \to 3m-2 \to 3m-1.
\]
So the exceptional actual lift has a forced landing row and a forced interface
continuation.
\end{proof}

\begin{corollary}[single-row globalization criterion]
Assume the standing inputs of Section~2. Then the full raw globalization
theorem is equivalent to one statement:
\begin{quote}
every actual lift of the exceptional cutoff row $\delta=3m-3$ glues through
$3m-2 \to 3m-1$ into the regular continuing endpoint class.
\end{quote}
\end{corollary}

\begin{proof}
By Theorem~5.2, the regular sector is already closed. By Input~E, every
exceptional full chain before the cutoff continues internally by
$\delta\mapsto \delta+1$. So the only row that can still launch a new endpoint
sheet is the exceptional cutoff row.
\end{proof}

\section{Final globalization theorem}

\begin{lemma}[one-chain pullback of endpoint class]
Let $C_{1,+}$ and $C_{2,+}$ be actual lifts of the same realized boundary label
$\delta+1$, and suppose they lie in the same endpoint class.

Assume $C_1$ and $C_2$ are predecessor full chains of $C_{1,+}$ and $C_{2,+}$
respectively, both with boundary label $\delta$.

Then $C_1$ and $C_2$ also lie in the same endpoint class.
\end{lemma}

\begin{proof}
By Lemma~4.2, the current-event word emitted along any full chain of boundary
label $\delta$ is the same canonical block $W(\delta)$. So the padded future
word of each predecessor chain is obtained by prefixing the same block
$W(\delta)$ to the padded future word of its successor chain.

Since $C_{1,+}$ and $C_{2,+}$ lie in the same endpoint class, they have the
same padded future current-event word. Prefixing the same block $W(\delta)$ to
both words preserves equality. Therefore $C_1$ and $C_2$ have the same padded
future current-event word, hence lie in the same endpoint class.
\end{proof}

\begin{theorem}[odd-$m$ D5 final globalization theorem]
For odd $m$ in the accepted D5 regime, every actual lift of the exceptional
cutoff row $\delta=3m-3$ continues through
\[
3m-2 \to 3m-1
\]
and lies in the regular continuing endpoint class there.

Consequently:
\begin{enumerate}[label=\textup{(\arabic*)},leftmargin=2.5em]
\item no mixed-status realized $\delta$ exists;
\item fixed realized $\delta$ determines tail length;
\item fixed realized $\delta$ determines endpoint class and padded future word;
\item $\rho=\rho(\delta)$ globally;
\item raw global $(\beta,\delta)$ is exact on the true accessible boundary
      union.
\end{enumerate}
\end{theorem}

\begin{proof}
Take any accessible full chain.

If it is regular, Theorem~5.2 gives a forward successor.

If it is exceptional but not the cutoff chain, Input~E gives a forward
successor by $\delta\mapsto \delta+1$.

If it is the exceptional cutoff chain, Theorem~6.1 gives a forward successor
and pins its next two labels to $3m-2\to 3m-1$.

So every accessible full chain has a forward successor.

Now let $C$ be a splice-connected accessible component. By Input~C, its
realized boundary image is either a forward interval in $\Z/m^2\Z$ or the full
cycle. A finite forward interval has a right-end chain with no successor, which
we have just ruled out. Hence every splice-connected accessible component is
total, and every actual lift has tail length $\infty$.

Theorem~4.4 says that fixed realized $\delta$ determines the padded future
current-event word, the endpoint class, and the boundary right-congruence class
$\rho$ once tail length is fixed. Since tail length is now always $\infty$, it
follows that
\[
\rho=\rho(\delta)
\]
globally, and raw global $(\beta,\delta)$ is exact.

It remains to recover the stated gluing conclusion. Let $E_+$ be the exceptional
lift at $\delta=3m-1$ reached from the cutoff row, and let $R_+$ be the regular
source-$1$ start lift at the same label. By Theorem~6.1, $E_+$ exists and is
reached through $3m-3\to 3m-2\to 3m-1$. Both $E_+$ and $R_+$ have infinite tail
length, so Theorem~4.4 forces them to have the same endpoint class.
Lemma~7.1 then pulls that equality back one full chain to the interface row
$3m-2$. Thus every actual lift of the exceptional cutoff row glues through the
interface
$3m-2\to 3m-1$ into the regular continuing endpoint class.
\end{proof}

\section{Honest boundary}

\begin{remark}
This note is meant to be read in one pass, but it still has a sharp honesty
boundary.

What is fully carried out here:
\begin{enumerate}[label=\textup{(\arabic*)},leftmargin=2.5em]
\item the structural first-exit theorem from the explicit trigger formula and
      the coarse patch facts;
\item the fixed-$\delta$ tail-length reduction from the componentwise bridge
      law;
\item the chart/interface landing theorem;
\item the regular closure theorem;
\item the final globalization proof.
\end{enumerate}

What is stated here as standing theorem-package input rather than redeveloped
from raw artifacts:
\begin{enumerate}[label=\textup{(\arabic*)},leftmargin=2.5em]
\item the exact trigger-family formula itself;
\item the coarse patch-table facts;
\item the componentwise concrete bridge theorem;
\item the regular chart / endpoint support package;
\item the exceptional source-$3$ interior continuation fact.
\end{enumerate}

So this is a one-document globalization package for reading and manuscript
flow, not a first-principles rebuild of every older support artifact and not
yet a full graph-theoretic Hamilton-decomposition proof.
\end{remark}

\end{document}
